% Part: set-theory
% Chapter: ordinals
% Section: basic

\documentclass[../../../include/open-logic-section]{subfiles}

\begin{document}

\olfileid{sth}{ordinals}{basic}
\olsection{ক্রমসংখ্যার মৌলিক ধর্ম}

আমরা দেখেছি, প্রথম কয়েকটি ক্রমসংখ্যাই স্বাভাবিক
সংখ্যা। ক্রমসংখ্যার তত্ত্ব গড়ার প্রধান কারণ হলো
স্বাভাবিক সংখ্যার আবেশনীতি আরও বিস্তৃত করা।
কয়েকটি প্রাথমিক ফলের মধ্য দিয়ে সেখানে পৌঁছোব।

\begin{lem}\ollabel{ordmemberord}
ক্রমসংখ্যার প্রত্যেক !!{element} একটি ক্রমসংখ্যা।
\end{lem}

\begin{proof}
ধরি, $\alpha$ একটি ক্রমসংখ্যা এবং $b \in \alpha$।
$\alpha$ পরিযায়ী বলে $b \subseteq \alpha$।
তাই $\in$ যেমন $\alpha$-কে সুক্রমে বিন্যস্ত করে,
তেমনি $b$-কেও $\in$ সুক্রমে বিন্যস্ত করে।

$b$ যে পরিযায়ী, তা দেখতে ধরি $x \in c \in b$।
$b \subseteq \alpha$ বলে $c \in \alpha$। আবার
$\alpha$ পরিযায়ী বলে $c \subseteq \alpha$;
তাই $x \in \alpha$। অতএব $x, c, b \in \alpha$।
কিন্তু $\in$ $\alpha$-কে সুক্রমে বিন্যস্ত করে;
তাই \olref[sth][ordinals][wo]{wo:strictorder} অনুযায়ী
$\in$ হলো $\alpha$-এর উপর একটি পরিযায়ী সম্পর্ক।
সুতরাং $x \in c \in b$ থেকে $x \in b$।
সার্বিকভাবে, $c \subseteq b$।
\end{proof}

\begin{cor}\ollabel{ordissetofsmallerord}
যেকোনো ক্রমসংখ্যা~$\alpha$-এর জন্য
$\alpha = \Setabs{\beta \in \alpha}{\beta \text{ একটি ক্রমসংখ্যা}}$।
\end{cor}

\begin{proof}
\olref{ordmemberord} থেকে সরাসরি।
\end{proof}

পরের দুটি প্রধান ফল,
\olref{ordinductionschema} এবং \olref{ordtrichotomy},
মোটামুটি বলে যে সদস্যতা সম্পর্কে ক্রমসংখ্যাগুলি
নিজেরাই সুক্রমিত:

\begin{thm}[ট্রান্সফাইনাইট আবেশ]\ollabel{ordinductionschema}
যেকোনো সূত্র $\phi(x)$-এর জন্য:
\[
	\text{যদি }\exists \alpha \phi(\alpha)\text{, তবে }\exists \alpha(\phi(\alpha)
	\land  (\forall \beta \in \alpha) \lnot \phi(\beta))
\]
এখানে প্রদর্শিত পরিমাণসূচকগুলি অপ্রকাশিতভাবে
ক্রমসংখ্যায় সীমাবদ্ধ।
\end{thm}

\begin{proof}
%\olref{ordinductionschema1}.
ধরি, কোনো ক্রমসংখ্যা $\alpha$-এর জন্য
$\phi(\alpha)$ সত্য। যদি $ (\forall \beta
\in \alpha) \lnot \phi(\beta)$ হয়, কাজ শেষ।
নতুবা $\alpha$ ক্রমসংখ্যা বলে তার মধ্যে
% BN-SRC-436: source says a least element “which is phi”; the intended condition is phi(beta).
$\phi$-সত্য এমন সদস্যদের একটি $\in$-ক্ষুদ্রতম
!!{element} আছে; \olref{ordmemberord} অনুযায়ী সেটিও
ক্রমসংখ্যা।
%	\olref{ordinductionschema2}। ধরি কোনো ক্রমসংখ্যা
%	$\gamma$ আছে যাতে $\lnot\phi(\gamma)$। তাহলে
%	\olref{ordinductionschema1} থেকে এমন একটি $\in$-ন্যূনতম
%	$\alpha$ আছে যাতে $\lnot\phi(\alpha)$। তাই $(\forall \beta <
%	\alpha) \phi(\beta)$; ফলে শর্তবাক্যের পূর্বপক্ষ
%	false.
\end{proof}
\noindent
লক্ষ করো, \olref{ordinductionschema}-কে সমানভাবে
নিচের ছক দিয়েও প্রকাশ করা যায়:
\[
\text{যদি }\forall \alpha((\forall \beta \in \alpha)\phi(\beta) \lif
\phi(\alpha))\text{, তবে }\forall \alpha\phi(\alpha)
\]
\olref{ordinductionschema}-এ $\lnot\phi(\alpha)$
বসিয়ে তারপর প্রাথমিক যৌক্তিক রূপান্তর করলেই এটি
পাওয়া যায়।

\begin{thm}[ত্রিবিভাজন]\ollabel{ordtrichotomy}
যেকোনো ক্রমসংখ্যা $\alpha$ ও $\beta$-এর জন্য
$\alpha \in \beta \lor \alpha = \beta \lor \beta \in \alpha$।
\end{thm}

\begin{proof}
প্রমাণে দ্বিস্তর আবেশ, অর্থাৎ
\olref{ordinductionschema} দুবার ব্যবহার করব।
$x$-কে $y$-এর সঙ্গে \emph{তুলনীয়} বলব যদি
$x \in y \lor x = y \lor y \in x$।

আবেশের জন্য ধরি, $\alpha$-এর প্রতিটি ক্রমসংখ্যা
\emph{প্রত্যেক} ক্রমসংখ্যার সঙ্গে তুলনীয়। আরও একবার
আবেশের জন্য ধরি, $\alpha$ হলো $\beta$-এর প্রতিটি
ক্রমসংখ্যার সঙ্গে তুলনীয়। আমরা দেখাব $\alpha$
ও $\beta$ তুলনীয়। তাহলে $\beta$-র উপর আবেশে
পাব, $\alpha$ প্রত্যেক ক্রমসংখ্যার সঙ্গে তুলনীয়;
তারপর $\alpha$-র উপর আবেশে পাব, \emph{প্রত্যেক}
ক্রমসংখ্যা \emph{প্রত্যেক} ক্রমসংখ্যার সঙ্গে তুলনীয়।
এটাই চাওয়া ফল। এখন $\alpha \notin \beta$ ও
$\beta \notin \alpha$ ধরে $\alpha = \beta$ দেখালেই
যথেষ্ট।

$\alpha \subseteq \beta$ দেখাতে একটি
$\gamma \in \alpha$ নিই। \olref{ordmemberord}
অনুযায়ী এটিও ক্রমসংখ্যা। তাই প্রথম আবেশ-অনুমানে
$\gamma$ ও $\beta$ তুলনীয়। কিন্তু $\gamma
= \beta$ অথবা $\beta \in \gamma$ হলে
$\beta \in \alpha$ (দরকারে $\alpha$-র
পরিযায়িতা ব্যবহার করে)—যা আমাদের অনুমানের
বিরোধী। তাই $\gamma \in \beta$। সার্বিকভাবে,
$\alpha \subseteq \beta$।

দ্বিতীয় আবেশ-অনুমান ব্যবহার করে একেবারে
অনুরূপ যুক্তিতে $\beta \subseteq \alpha$ পাই।
তাই $\alpha = \beta$।
\end{proof}\noindent
এই জন্য কখনো $\alpha \in \beta$-র বদলে
$\alpha <\beta$ লিখব, কারণ এখানে $\in$ ক্রমসম্পর্কের
মতো কাজ করছে। এতে কোনো গভীর কারণ নেই; সংকেতটি
পরিচিত, আর $\alpha \in \beta \lor \alpha = \beta$-র
চেয়ে $\alpha \leq \beta$ লেখা সহজ।\footnote{চাইলে
$\alpha \mathrel{\underline{\in}} \beta$ লিখতে
পারতাম; কিন্তু তা একেবারেই অপ্রচলিত।}

এই ফলগুলির দ্রুত পাওয়া দুটি পরিণাম দেখি;
প্রথমটিতে নতুন সংকেতলিপিটি কাজে লাগবে:

\begin{cor}\ollabel{ordordered}
যদি $\exists \alpha\phi(\alpha)$ হয়, তবে
$\exists \alpha(\phi(\alpha)
\land \forall \beta(\phi(\beta) \lif \alpha \leq \beta))$।
আর যেকোনো ক্রমসংখ্যা $\alpha, \beta, \gamma$-এর
জন্য $\alpha \notin \alpha$ এবং
$\alpha \in \beta \in \gamma \lif \alpha \in \gamma$
উভয়ই সত্য।
\end{cor}

\begin{proof}
\olref[wo]{wo:strictorder}-এর মতোই।
\end{proof}

\begin{prob}
\olref[sth][ordinals][basic]{ordtrichotomy}-এর প্রমাণে
``একেবারে অনুরূপ যুক্তি'' অংশটি সম্পূর্ণ করো।
\end{prob}

\begin{cor}\ollabel{corordtransitiveord}
$A$ একটি ক্রমসংখ্যা যদি এবং কেবল যদি $A$
ক্রমসংখ্যাগুলির একটি পরিযায়ী সেট।
\end{cor}

\begin{proof}
\emph{বাঁ থেকে ডানে।} \olref{ordmemberord}
অনুযায়ী। \emph{ডান থেকে বাঁয়ে।} $A$ যদি
ক্রমসংখ্যাগুলির একটি পরিযায়ী সেট হয়, তবে
\olref{ordinductionschema} এবং
\olref{ordtrichotomy} অনুযায়ী $\in$ $A$-কে
সুক্রমে বিন্যস্ত করে।
\end{proof}

আমরা \olref{ordinductionschema} ও
\olref{ordtrichotomy}-এর সারমর্ম বলেছি:
$\in$ ক্রমসংখ্যাগুলিকে সুক্রমে বিন্যস্ত করে।
কিন্তু নিচের ফলের জন্য এ রকম কথা বলার সময়
\emph{বিশেষ সতর্ক} থাকতে হবে:

\begin{thm}[বুরালি-ফোর্তির কূটাভাস]\ollabel{buraliforti}
সব ক্রমসংখ্যার কোনো সেট নেই।
\end{thm}

\begin{proof}
বিরোধের উদ্দেশ্যে ধরি, $O$ সব ক্রমসংখ্যার সেট।
যদি $\alpha \in \beta \in O$ হয়, তবে
\olref{ordmemberord} অনুযায়ী $\alpha$ একটি
ক্রমসংখ্যা; তাই $\alpha \in O$। অতএব $O$
পরিযায়ী; তাই \olref{corordtransitiveord} অনুযায়ী
$O$ একটি ক্রমসংখ্যা। ফলে $O \in O$,
যা \olref{ordordered}-এর বিরোধী।
\end{proof}
\noindent
ফলটি \citeauthor{Burali-Forti1897}-র নামে।
কিন্তু ১৮৯৯ সালে ডেডেকিন্ডকে লেখা একটি চিঠিতে
কান্টরই প্রথম সব ক্রমসংখ্যার সেট ধরে নেওয়ার
\emph{বিরোধটি} স্পষ্টভাবে দেখেন। ফন হেইয়েনুর্টের
ব্যাখ্যায়:
\begin{quote}
  বুরালি-ফোর্তি নিজে মনে করেছিলেন,
  \emph{অসঙ্গতি দেখিয়ে} এই সিদ্ধান্ত প্রতিষ্ঠিত হয় যে
  ক্রমসংখ্যাগুলির স্বাভাবিক ক্রম কেবল একটি আংশিক ক্রম।
  \citep[p.~105]{Heijenoort1967}
\end{quote}
বুরালি-ফোর্তির ভুলটি সরিয়ে রেখে আগের আলোচনা
সংক্ষেপে বলি: প্রত্যেক ক্রমসংখ্যা একটি সেট, যা
সদস্যতা সম্পর্কে নিজে সুক্রমিত; ক্রমসংখ্যাগুলি
সমষ্টিগতভাবেও সদস্যতা সম্পর্কে সুক্রমিত, কিন্তু
সব মিলে কোনো সেট গঠন করে না।

সবশেষে \olref{ordinductionschema} ও
\olref{ordtrichotomy} থেকে ক্রমসংখ্যার আরও
কয়েকটি মৌলিক ধর্ম দেখি।

\begin{prop}
ক্রমসংখ্যাগুলির যেকোনো কঠোরভাবে নিম্নগামী
অনুক্রম সসীম।
\end{prop}

\begin{proof}
ক্রমসংখ্যাগুলির কোনো অসীম কঠোরভাবে নিম্নগামী
অনুক্রম $\alpha_0 > \alpha_1 > \alpha_2 > \ldots$-এর
কোনো $<$-ন্যূনতম সদস্য নেই; এটি
\olref{ordinductionschema}-এর বিরোধী।
\end{proof}

\begin{prop}\ollabel{ordinalsaresubsets}
যেকোনো ক্রমসংখ্যা $\alpha, \beta$-র জন্য
$\alpha \subseteq \beta \lor \beta \subseteq \alpha$।
\end{prop}

\begin{proof}
যদি $\alpha \in \beta$ হয়, তবে $\beta$
পরিযায়ী বলে $\alpha \subseteq \beta$।
একইভাবে, $\beta \in \alpha$ হলে
$\beta \subseteq \alpha$। আর $\alpha = \beta$
হলে $\alpha \subseteq \beta$ ও
$\beta \subseteq \alpha$। তাই
\olref{ordtrichotomy} থেকে ফলটি পাই।
\end{proof}

\begin{prop}\ollabel{ordisoidentity}
যেকোনো ক্রমসংখ্যা $\alpha, \beta$-র জন্য
$\alpha = \beta$ যদি এবং কেবল যদি
$\ordeq{\alpha}{\beta}$।
\end{prop}

\begin{proof}
ক্রমসংখ্যাগুলি সুক্রমবিন্যাস; তাই
ত্রিবিভাজন (\olref[basic]{ordtrichotomy}) এবং
\olref[iso]{wellordnotinitial} থেকে সরাসরি।
\end{proof}

\begin{prob}\ollabel{probunionordinalsordinal}
প্রমাণ করো: $X$-এর প্রত্যেক সদস্য ক্রমসংখ্যা
হলে $\bigcup X$-ও ক্রমসংখ্যা।
\end{prob}

\end{document}
